\color{unidarkgreen}{\subsection{Cигнализация}}
\color{black}

\subsubsection{Общие сведения}

\begin{enumerate}[label=\arabic{section}.\arabic{subsection}.\arabic{subsubsection}.\arabic*,  align=left, labelsep=4pt, leftmargin=0pt, itemindent=57pt]

\item
Устройство формирует выходной сигнал на сигнализацию по логической схеме «ИЛИ» при
наступлении указанных ниже событий.

\item
Сигналы «Неисправность» и «Неисправност.фикс.» формируются при возникновении одного из
следующих событий:
\begin{itemize}
\item Срабатывание ГЗ;
\item Срабатывание ТЗ;
\item Срабатывание КЦТ;
\item Срабатывание КЦН;
\item Срабатывание КИ;
\item Срабатывание КОН;
\item Неисправности ОТ;
\item Срабатывание Защит ЭМО и ЭМВ;
\item Неисправности ЦУ;
\item Несоответствии ПУ;
\item Потеря GOOSE.
\end{itemize}

\item
Сигналы «Срабатывание» и «срабатывание имп.» формируются при срабатывании одной из
следующих защит или функций:
\begin{itemize}
\item Срабатывание ДТЗ;
\item Срабатывание МТЗ;
\item Срабатывание ГЗ;
\item Срабатывание ТЗ;
\item Срабатывание ЗНР;
\item Срабатывание УРОВ;
\item Срабатывание ТЗНП;
\item Срабатывание ЗПО;
\item Срабатывание ЗПТ;
\item Срабатывание РТПО;
\item Отключение В;
\item Отключение смежного Т.
\end{itemize}

\begin{figure}[H]
\centering
\includegraphics[width=0.8\textwidth,height=0.8\textheight,keepaspectratio]{\fbpath/91 Сигнализация/img/SIGN.pdf}
\captionof{figure}{Функциональная схема предупредительной сигнализации}\label{pic:SIGN}
\end{figure}

\end{enumerate}